\chapter{1958年云南卷}

\begin{problem}
解方程 \(\left(\log_{x}\sqrt{5}\right)^{2}+\log_{x}5-\frac{5}{4}=0\).
\end{problem}

\begin{solution}
由对数的定义，须有 \(x>0\) 且 \(x\ne1\). 设 \(t=\log_{x}5\)，则
\(\log_{x}\sqrt{5}=\frac{1}{2}\log_{x}5=\frac{t}{2}\). 原方程等价于
\[
\frac{t^{2}}{4}+t-\frac{5}{4}=0
\]
即 \(t^{2}+4t-5=0\)，所以
\[
(t-1)(t+5)=0
\]

当 \(t=1\) 时，\(\log_{x}5=1\)，故 \(x=5\).

当 \(t=-5\) 时，\(\log_{x}5=-5\)，故 \(x^{-5}=5\)，从而 \(x^{5}=\frac{1}{5}\)，所以 \(x=\sqrt[5]{\frac{1}{5}}=\frac{\sqrt[5]{625}}{5}\). 这两个值均满足 \(x>0\) 且 \(x\ne1\)，故原方程的解为
\(x=5\) 或 \(x=\frac{\sqrt[5]{625}}{5}\).
\end{solution}

\begin{problem}
求 \((1+x)+(1+x)^{2}+\cdots +(1+x)^{10}\) 式中 \(x^{8}\) 项的系数（不必求出其它各项的系数）.
\end{problem}

\begin{solution}
在 \((1+x)^{n}\) 中，二项式定理给出的 \(x^{8}\) 项系数为 \(\mathrm C_{n}^{8}\)，但只有 \(n\ge8\) 时该项才存在. 因此原式中只有 \((1+x)^{8}\)、\((1+x)^{9}\)、\((1+x)^{10}\) 含有 \(x^{8}\) 项.

其中，\((1+x)^{8}\) 中 \(x^{8}\) 的系数为 \(\mathrm C_{8}^{8}=1\)；\((1+x)^{9}\) 中 \(x^{8}\) 的系数为 \(\mathrm C_{9}^{8}=\mathrm C_{9}^{1}=9\)；\((1+x)^{10}\) 中 \(x^{8}\) 的系数为
\[
\mathrm C_{10}^{8}=\mathrm C_{10}^{2}=\frac{10\times9}{1\times2}=45
\]
所以所求系数为 \(1+9+45=55\).
\end{solution}

\begin{problem}
求圆锥曲线 \(2x^{2}-8xy-4y^{2}-4y+1=0\) 之焦点及准线.
\end{problem}

\begin{solution}
先平移坐标轴. 令
\(x=x'+a\)，\(y=y'+b\).
代入原方程并按 \(x'\)、\(y'\) 的次数整理，得
\[
2x'^{2}-8x'y'-4y'^{2}+(4a-8b)x'-(8a+8b+4)y'+2a^{2}-8ab-4b^{2}-4b+1=0
\]
令一次项系数为零，得到
\[
\begin{cases}
4a-8b=0,\\
8a+8b+4=0.
\end{cases}
\]
解得 \(a=-\frac{1}{3}\)，\(b=-\frac{1}{6}\). 因而
\(x'=x+\frac{1}{3}\)，\(y'=y+\frac{1}{6}\)，且方程化为
\[
3x'^{2}-12x'y'-6y'^{2}+2=0
\]

再旋转坐标轴. 对二次项 \(3x'^{2}-12x'y'-6y'^{2}\)，有
\(A=3\)，\(B=-12\)，\(C=-6\)，故取 \(\theta\) 使
\(\tan2\theta=\frac{B}{A-C}=-\frac{4}{3}\)，\(\cos2\theta=-\frac{3}{5}\)，\(\sin2\theta=\frac{4}{5}\).
取第一象限的 \(\theta\)，则
\(\cos\theta=\frac{1}{\sqrt5}\)，\(\sin\theta=\frac{2}{\sqrt5}\).
令
\(x'=x''\cos\theta-y''\sin\theta\)，\(y'=x''\sin\theta+y''\cos\theta\)，
即
\(x'=\frac{x''-2y''}{\sqrt5}\)，\(y'=\frac{2x''+y''}{\sqrt5}\).
代入平移后的方程，化简得
\[
9x''^{2}-6y''^{2}=2
\]
即
\[
\frac{x''^{2}}{2/9}-\frac{y''^{2}}{1/3}=1
\]
\solutionfigure{\bitmapfigure[width=0.22\linewidth]{img/1958/yunnan/q3_solution_fig1.png}}

这是以原点为中心、实半轴平方 \(a_{0}^{2}=\frac{2}{9}\)、虚半轴平方 \(b_{0}^{2}=\frac{1}{3}\) 的双曲线. 于是焦距参数满足
\[
c_{0}=\sqrt{a_{0}^{2}+b_{0}^{2}}=\sqrt{\frac{2}{9}+\frac{1}{3}}=\frac{\sqrt5}{3}
\]
在 \(x''Oy''\) 坐标系中，两个焦点为
\(F_{1}\left(\frac{\sqrt5}{3},0\right)\)，\(F_{2}\left(-\frac{\sqrt5}{3},0\right)\).
由坐标变换，在焦点处 \(y''=0\)，且
\(x'=\frac{x''}{\sqrt5}\)，\(y'=\frac{2x''}{\sqrt5}\).
因此，当 \(x''=\frac{\sqrt5}{3}\) 时，\((x',y')=\left(\frac{1}{3},\frac{2}{3}\right)\)，从而
\((x,y)=\left(0,\frac{1}{2}\right)\)；当 \(x''=-\frac{\sqrt5}{3}\) 时，\((x',y')=\left(-\frac{1}{3},-\frac{2}{3}\right)\)，从而
\((x,y)=\left(-\frac{2}{3},-\frac{5}{6}\right)\).

又因为
\[
x''=\frac{x'+2y'}{\sqrt5}=\frac{x+2y+2/3}{\sqrt5}
\]
双曲线在旋转坐标系中的两条准线
\(x''=\pm\frac{a_{0}^{2}}{c_{0}}=\pm\frac{2}{3\sqrt5}\) 分别给出
\(x+2y+\frac{2}{3}=\frac{2}{3}\)，\(x+2y+\frac{2}{3}=-\frac{2}{3}\).
所以原坐标系中的焦点为
\(\left(0,\frac{1}{2}\right)\)，\(\left(-\frac{2}{3},-\frac{5}{6}\right)\)，
准线为
\(x+2y=0\)，\(x+2y+\frac{4}{3}=0\).
\solutionfigure{\bitmapfigure[width=0.52\linewidth]{img/1958/yunnan/q3_solution_fig2.png}}
\end{solution}

\begin{problem}
求圆锥曲线 \(x^{2}+y^{2}=49\) 及 \(x^{2}+y^{2}-20y+91=0\) 之公切线.
\end{problem}

\begin{solution}
两圆分别为
\(\mathcal C_{1}:x^{2}+y^{2}=7^{2}\)，\(\mathcal C_{2}:x^{2}+(y-10)^{2}=3^{2}\).
因此它们的圆心和半径分别为 \(O_{1}(0,0)\)，\(r_{1}=7\) 及 \(O_{2}(0,10)\)，\(r_{2}=3\).
\solutionfigure{\bitmapfigure[width=0.28\linewidth]{img/1958/yunnan/q4_solution_fig1.png}}

先说明所求公切线不是竖直线. 若直线为 \(x=d\)，它与第一圆相切需有 \(|d|=7\)，而与第二圆相切需有 \(|d|=3\)，两者不能同时成立. 因此设公切线为 \(y=kx+b\)，即
\(kx-y+b=0\).

由点到直线的距离等于相应半径，得
\(\frac{|b|}{\sqrt{k^{2}+1}}=7\)，\(\frac{|b-10|}{\sqrt{k^{2}+1}}=3\).
两式平方后分别为
\(b^{2}=49(k^{2}+1)\)，\((b-10)^{2}=9(k^{2}+1)\).
消去 \(k^{2}+1\)，得
\[
49(b-10)^{2}=9b^{2}
\]
即
\(40b^{2}-980b+4900=0\)，化简为 \(2b^{2}-49b+245=0\).
所以
\[
b=\frac{49\pm\sqrt{49^{2}-4\times2\times245}}{4}=\frac{49\pm21}{4}
\]
即 \(b=\frac{35}{2}\) 或 \(b=7\).

当 \(b=\frac{35}{2}\) 时，代入 \(b^{2}=49(k^{2}+1)\)，得
\(k^{2}+1=\frac{(35/2)^{2}}{49}=\frac{25}{4}\)，所以 \(k=\pm\frac{\sqrt{21}}{2}\).
于是得到两条切线
\(y=\frac{\sqrt{21}}{2}x+\frac{35}{2}\)，\(y=-\frac{\sqrt{21}}{2}x+\frac{35}{2}\)，
即
\(\sqrt{21}x-2y+35=0\)，\(\sqrt{21}x+2y-35=0\).

当 \(b=7\) 时，\(49=49(k^{2}+1)\)，故 \(k=0\)，得到切线 \(y=7\). 代回距离条件可知这三条直线均确为公切线，其中 \(y=7\) 是两圆外切点处的内公切线，另外两条是外公切线. 因此所求公切线为
\(y=7\)，\(\sqrt{21}x-2y+35=0\)，\(\sqrt{21}x+2y-35=0\).
\end{solution}

\begin{problem}
设一四面体为一平面所截，若截面是平行四边形，则此截面平行于四面体的两条棱.
\end{problem}

\begin{solution}
设四面体为 \(ABCD\)，截面为平行四边形 \(EFGH\)，其中 \(E\)、\(F\)、\(G\)、\(H\) 分别在 \(AB\)、\(AC\)、\(CD\)、\(BD\) 上，如图所示.
\solutionfigure{\bitmapfigure[width=0.30\linewidth]{img/1958/yunnan/q5_solution_fig1.png}}

因为 \(EFGH\) 是平行四边形，所以 \(EF\parallel HG\). 直线 \(HG\) 在平面 \(BCD\) 内，直线 \(EF\) 在平面 \(ABC\) 内，且平面 \(ABC\) 与平面 \(BCD\) 的交线为 \(BC\). 由直线与平面的平行性质，得到
\[
EF\parallel BC
\]
直线 \(EF\) 在截面平面 \(EFGH\) 内，因此截面平面平行于棱 \(BC\).

平行四边形的另一组对边满足 \(FG\parallel EH\)，其中 \(FG\) 在平面 \(ACD\) 内，\(EH\) 在平面 \(ABD\) 内，而两平面的交线为 \(AD\)，所以
\[
FG\parallel AD
\]
直线 \(FG\) 在截面平面 \(EFGH\) 内，因此截面平面也平行于棱 \(AD\). 综上，截面平行于四面体的两条棱 \(BC\) 和 \(AD\).
\end{solution}

\begin{problem}
设三角形 \(ABC\) 中，\(\angle A\) 的外角平分线与 \(BC\) 边之延长线相交于 \(D\)，与外接圆周相交于 \(E\)，则 \(AB\cdot AC=AE\cdot AD\).
\end{problem}

\begin{solution}
连接 \(BE\)，并延长 \(AB\) 到 \(F\). 由于 \(AD\) 是 \(\angle CAF\) 的平分线，且 \(A\)，\(B\)，\(F\) 共线，所以
\[
\angle BAE=\angle DAF=\angle CAD
\]
又因为四边形 \(ACBE\) 内接于圆，所以
\(\angle AEB+\angle ACB=180^\circ\). 由于 \(B\)，\(C\)，\(D\) 共线，\(\angle ACB+\angle ACD=180^\circ\)，故
\[
\angle AEB=\angle ACD
\]
于是 \(\triangle ABE\sim\triangle ADC\)，对应关系为 \(A\leftrightarrow A\)，\(B\leftrightarrow D\)，\(E\leftrightarrow C\). 因而
\[
\frac{AB}{AE}=\frac{AD}{AC}
\]
交叉相乘即得
\[
AB\cdot AC=AE\cdot AD
\]
\end{solution}

\begin{problem}
设有互相外切的二圆的半径各为 \(a\) 与 \(b\)（\(a>b\)），作此二圆的两外公切线，设其夹角为 \(A\)，证明：\(\sin A=\frac{4(a-b)\sqrt{ab}}{(a+b)^{2}}\).
\end{problem}

\begin{solution}
设两圆圆心为 \(O_{1}\)、\(O_{2}\)，则
\[
O_{1}O_{2}=a+b
\]
取一条外公切线，分别从 \(O_{1}\)、\(O_{2}\) 向它作垂线，垂足为 \(B\)、\(C\). 因为两圆心在公切线同侧，且 \(O_{1}B=a\)、\(O_{2}C=b\)，所以线段 \(O_{1}O_{2}\) 在垂直于公切线方向上的投影长度为 \(a-b\). 两条外公切线的交点与两个圆心共线，故圆心连线平分两条外公切线的夹角. 于是圆心连线与该外公切线的夹角为 \(\frac{A}{2}\)，从而
\[
\sin\frac{A}{2}=\frac{a-b}{a+b}
\]
又由平方关系
\[
\cos\frac{A}{2}=\sqrt{1-\sin^{2}\frac{A}{2}}=\frac{\sqrt{(a+b)^{2}-(a-b)^{2}}}{a+b}=\frac{2\sqrt{ab}}{a+b}
\]
因此
\[
\sin A=2\sin\frac{A}{2}\cos\frac{A}{2}=2\times\frac{a-b}{a+b}\times\frac{2\sqrt{ab}}{a+b}=\frac{4(a-b)\sqrt{ab}}{(a+b)^{2}}
\]
\solutionfigure{\bitmapfigure[width=0.52\linewidth]{img/1958/yunnan/q7_solution_fig1.png}}
\end{solution}

\begin{problem}
在三角形 \(ABC\) 中，\(r\) 与 \(R\) 分别为其内切圆与外接圆的半径，求证：\(r=4R\sin\frac{A}{2}\sin\frac{B}{2}\sin\frac{C}{2}\).
\end{problem}

\begin{solution}
设内切圆圆心为 \(O\)，它与 \(AB\) 相切于 \(D\). 连接 \(OA\)、\(OB\)、\(OD\)，则 \(OD\perp AB\)，且 \(OD=r\). 由内心在角平分线上，\(\angle OAD=\frac{A}{2}\)，\(\angle OBD=\frac{B}{2}\).
\solutionfigure{\bitmapfigure[width=0.28\linewidth]{img/1958/yunnan/q8_solution_fig1.png}}

在直角三角形 \(AOD\) 中，
\[
\cot\frac{A}{2}=\frac{AD}{OD}
\]
所以 \(AD=r\cot\frac{A}{2}\). 在直角三角形 \(BOD\) 中，
\[
\cot\frac{B}{2}=\frac{BD}{OD}
\]
故 \(BD=r\cot\frac{B}{2}\). 因而以 \(c=AB\) 表示边 \(AB\)，有
\[
c=AD+BD=r\left(\cot\frac{A}{2}+\cot\frac{B}{2}\right)
\]
将余切写成正弦、余弦，得
\[
c=r\frac{\sin\frac{B}{2}\cos\frac{A}{2}+\sin\frac{A}{2}\cos\frac{B}{2}}{\sin\frac{A}{2}\sin\frac{B}{2}}=r\frac{\sin\frac{A+B}{2}}{\sin\frac{A}{2}\sin\frac{B}{2}}
\]
因为 \(A+B+C=180^\circ\)，所以 \(\frac{A+B}{2}=90^\circ-\frac{C}{2}\)，从而
\[
c=r\frac{\cos\frac{C}{2}}{\sin\frac{A}{2}\sin\frac{B}{2}}
\]

由正弦定理，\(c=2R\sin C\). 再由二倍角公式，
\[
c=2R\sin C=4R\sin\frac{C}{2}\cos\frac{C}{2}
\]
两种表示相等，且三角形内角满足 \(0<C<180^\circ\)，所以 \(\cos\frac{C}{2}>0\)，可以约去该因子，得到
\[
4R\sin\frac{C}{2}=\frac{r}{\sin\frac{A}{2}\sin\frac{B}{2}}
\]
即
\[
r=4R\sin\frac{A}{2}\sin\frac{B}{2}\sin\frac{C}{2}
\]
\end{solution}
